% GENERATED FILE. Edit canonical lessons, then run npm run generate:books.
% canonical-source-hash: fnv1a64:5d83e5ef232fe059
% canonical-lessons: KA-C318-avalambisu, KA-C318-mele-avalambiside, KA-C318-iddiddare, KA-W318-six-lines, KA-R318-depends-and-would

\chapter{It Depends, and What Would Have Been}
\label{ch:ka-it-depends-and-what-would-have-been}
\hlchaptermodality{318}

\begin{chapteropening}
\textbf{By the end of this chapter:} \emph{I can say what something depends on with … mēle avalambiside, say what would have happened with -iddare … -uttidde, and write six connected lines about a trip.}

Complete the last lesson of chapter 318: say what something depends on and what would have been, after writing six connected lines about a trip.
\end{chapteropening}

\section[\texorpdfstring{\kn{ಅವಲಂಬಿಸು} (avalambisu)}{ಅವಲಂಬಿಸು (avalambisu)}]{\kn{ಅವಲಂಬಿಸು} (avalambisu) — to depend, to rely}

\label{lesson:KA-C318-avalambisu}



\begin{quote}
Before the new one: say \textquotedblleft{}on top, on\textquotedblright{} (\emph{mēle}) and \textquotedblleft{}according to\textquotedblright{} (\emph{prakāra}). Then \textquotedblleft{}If it rains tomorrow, I will stay at home.\textquotedblright{} (\emph{nāḷe maḷe bandare nānu maneyallē iruttēne}.)
\end{quote}

\subsection*{You'll want to know: \kn{ಅವಲಂಬಿಸು}}
\textbf{\kn{ಅವಲಂಬಿಸು}} — \emph{avalambisu} — \textquotedblleft{}to depend (on), to rely (on)\textquotedblright{}.

\begin{quote}
\textbf{\kn{ನಮ್ಮ} \kn{ಊರು} \kn{ಮಳೆಯ} \kn{ಮೇಲೆ} \kn{ಅವಲಂಬಿಸಿದೆ}.}
\end{quote}

\begin{quote}
\emph{namma ūru maḷeya mēle avalambiside.} — \textquotedblleft{}Our town depends on the rain.\textquotedblright{}
\end{quote}

What you depend on takes \textbf{\kn{ಮೇಲೆ}}, \textquotedblleft{}on\textquotedblright{}: \textbf{\kn{ಮಳೆಯ} \kn{ಮೇಲೆ}} (\emph{maḷeya mēle}, \textquotedblleft{}on the rain\textquotedblright{}). \textbf{\kn{ಅವಲಂಬಿಸಿದೆ}} is \textbf{\kn{ಅವಲಂಬಿಸಿ}}, the \textbf{-\kn{ಇ}} form, with \textbf{\kn{ಇದೆ}}: \textquotedblleft{}it is, having come to rest on\textquotedblright{}, so \textquotedblleft{}it depends\textquotedblright{}. Written the same, \textbf{\kn{ಅವಲಂಬಿಸಿದೆ}} is also \textquotedblleft{}I depended\textquotedblright{}; the sentence tells you which.

For \textquotedblleft{}it depended\textquotedblright{}, \textbf{\kn{ಇದೆ}} becomes \textbf{\kn{ಇತ್ತು}}:

\begin{quote}
\textbf{\kn{ಆ} \kn{ಹಳ್ಳಿ} \kn{ಒಂದು} \kn{ಕೆರೆಯ} \kn{ಮೇಲೆ} \kn{ಅವಲಂಬಿಸಿತ್ತು}.}
\end{quote}

\begin{quote}
\emph{ā haḷḷi ondu kereya mēle avalambisittu.} — \textquotedblleft{}That village depended on one tank.\textquotedblright{}
\end{quote}

\begin{cousinweb}
\textbf{\kn{ಅವಲಂಬಿಸು}} (\emph{avalambisu}) is Sanskrit \emph{avalamba} made into a Kannada verb with \textbf{-\kn{ಇಸು}}, as \textbf{\kn{ಯೋಚಿಸು}} is. CDIAL 826 gives \emph{avalamba} as \textquotedblleft{}hanging down\textquotedblright{} and, as a noun, \textquotedblleft{}stay, support\textquotedblright{}; Monier-Williams adds \textquotedblleft{}depending, resting upon\textquotedblright{}. To depend on something is to hang from it, which is also where English \emph{depend} comes from: Latin \emph{dēpendēre}, \textquotedblleft{}to hang down\textquotedblright{}.
\end{cousinweb}

\subsection*{Your turn}
Say these aloud:

\begin{itemize}
\raggedright
  \item \emph{avalambisu}, slowly, then at speed
  \item \textquotedblleft{}Our town depends on the rain\textquotedblright{} — \emph{namma ūru maḷeya mēle avalambiside}
  \item \textquotedblleft{}That village depended on one tank\textquotedblright{} — \emph{ā haḷḷi ondu kereya mēle avalambisittu}
\end{itemize}

\begin{tcolorbox}[breakable,colback=teal!4,colframe=teal!35!black,title={Before you move on}]
What does \textbf{\kn{ಅವಲಂಬಿಸು}} mean? (\textbf{To depend, to rely.}) What does the thing you depend on take? (\textbf{\kn{ಮೇಲೆ}}, \textquotedblleft{}on\textquotedblright{}.) \textquotedblleft{}It depends\textquotedblright{}? (\textbf{\kn{ಅವಲಂಬಿಸಿದೆ}}.) \textquotedblleft{}It depended\textquotedblright{}? (\textbf{\kn{ಅವಲಂಬಿಸಿತ್ತು}}.)
\end{tcolorbox}
\section[\texorpdfstring{\kn{ಮೇಲೆ} \kn{ಅವಲಂಬಿಸಿದೆ} (mēle avalambiside)}{ಮೇಲೆ ಅವಲಂಬಿಸಿದೆ (mēle avalambiside)}]{\kn{ಮೇಲೆ} \kn{ಅವಲಂಬಿಸಿದೆ} (mēle avalambiside) — it depends}

\label{lesson:KA-C318-mele-avalambiside}



\begin{quote}
Say \textquotedblleft{}to depend\textquotedblright{} (\emph{avalambisu}) and \textquotedblleft{}Our town depends on the rain.\textquotedblright{} (\emph{namma ūru maḷeya mēle avalambiside}.)
\end{quote}

\subsection*{You'll want to know: \kn{ಅದು} … \kn{ಮೇಲೆ} \kn{ಅವಲಂಬಿಸಿದೆ}}
\begin{quote}
\textbf{\kn{ನೀವು} \kn{ನಾಳೆ} \kn{ಬರುತ್ತೀರಾ}?}
\end{quote}

\begin{quote}
\emph{nīvu nāḷe baruttīrā?} — \textquotedblleft{}Will you come tomorrow?\textquotedblright{}
\end{quote}

\begin{quote}
\textbf{\kn{ಅದು} \kn{ಹವಾಮಾನದ} \kn{ಮೇಲೆ} \kn{ಅವಲಂಬಿಸಿದೆ}. \kn{ಮಳೆ} \kn{ಬಂದರೆ} \kn{ನಾನು} \kn{ಬರುವುದಿಲ್ಲ}.}
\end{quote}

\begin{quote}
\emph{adu havāmānada mēle avalambiside. maḷe bandare nānu baruvudilla.} — \textquotedblleft{}It depends on the weather. If it rains, I won't come.\textquotedblright{}
\end{quote}

\textbf{\kn{ಬರುತ್ತೀರಾ}} (\emph{baruttīrā}, \textquotedblleft{}will you come?\textquotedblright{}), said with respect; \textbf{\kn{ಬರುವುದಿಲ್ಲ}} (\emph{baruvudilla}, \textquotedblleft{}(I) won't come\textquotedblright{}). \textbf{\kn{ಅದು} … \kn{ಮೇಲೆ} \kn{ಅವಲಂಬಿಸಿದೆ}} is the reply when the honest answer is \textquotedblleft{}maybe\textquotedblright{}: name what it hangs on, then give the condition with \textbf{-\kn{ಅರೆ}}.

Shorter, and very common:

\begin{quote}
\textbf{\kn{ಅದು} \kn{ನಿಮ್ಮ} \kn{ಮೇಲೆ} \kn{ಅವಲಂಬಿಸಿದೆ}.}
\end{quote}

\begin{quote}
\emph{adu nimma mēle avalambiside.} — \textquotedblleft{}It's up to you.\textquotedblright{}
\end{quote}

Word for word, \textquotedblleft{}it depends on you\textquotedblright{}. Beside it sits \textbf{\kn{ಪ್ರಕಾರ}}, \textquotedblleft{}according to\textquotedblright{}, for what a plan or a rule says: \textbf{\kn{ವೇಳಾಪಟ್ಟಿಯ} \kn{ಪ್ರಕಾರ} \kn{ಬಸ್} \kn{ಆರು} \kn{ಗಂಟೆಗೆ} \kn{ಹೊರಡುತ್ತದೆ}} (\emph{vēḷāpaṭṭiya prakāra bas āru gaṇṭege horaḍuttade}, \textquotedblleft{}according to the timetable, the bus leaves at six\textquotedblright{}). Where \textbf{\kn{ಅವಲಂಬಿಸಿದೆ}} says \textquotedblleft{}it hangs on\textquotedblright{}, \textbf{\kn{ಪ್ರಕಾರ}} says \textquotedblleft{}it follows\textquotedblright{}.

\subsection*{Your turn}
Say these aloud:

\begin{itemize}
\raggedright
  \item \textquotedblleft{}Will you come tomorrow?\textquotedblright{} — \emph{nīvu nāḷe baruttīrā?}
  \item \textquotedblleft{}It depends on the weather. If it rains, I won't come.\textquotedblright{} — \emph{adu havāmānada mēle avalambiside. maḷe bandare nānu baruvudilla.}
  \item \textquotedblleft{}It's up to you\textquotedblright{} — \emph{adu nimma mēle avalambiside}
  \item one plan of yours that depends on something, and the condition
\end{itemize}

\begin{tcolorbox}[breakable,colback=teal!4,colframe=teal!35!black,title={Before you move on}]
\textquotedblleft{}It depends on the weather\textquotedblright{}? (\textbf{\kn{ಅದು} \kn{ಹವಾಮಾನದ} \kn{ಮೇಲೆ} \kn{ಅವಲಂಬಿಸಿದೆ}}.) \textquotedblleft{}It's up to you\textquotedblright{}? (\textbf{\kn{ಅದು} \kn{ನಿಮ್ಮ} \kn{ಮೇಲೆ} \kn{ಅವಲಂಬಿಸಿದೆ}}.) What comes after \textquotedblleft{}it depends\textquotedblright{}? (\textbf{The condition, with -\kn{ಅರೆ}.})
\end{tcolorbox}
\section[\texorpdfstring{\kn{ಇದ್ದಿದ್ದರೆ} … \kn{ಹೋಗುತ್ತಿದ್ದೆ} (iddiddare … …)}{ಇದ್ದಿದ್ದರೆ … ಹೋಗುತ್ತಿದ್ದೆ (iddiddare … …)}]{\kn{ಇದ್ದಿದ್ದರೆ} … \kn{ಹೋಗುತ್ತಿದ್ದೆ} (iddiddare … hōguttidde) — what would have been}

\label{lesson:KA-C318-iddiddare}



\begin{quote}
Say \textquotedblleft{}If it rains tomorrow, I will stay at home.\textquotedblright{} (\emph{nāḷe maḷe bandare nānu maneyallē iruttēne}.) Then \textquotedblleft{}I was reading\textquotedblright{}. (\emph{ōduttidde}.)

The first is a real \textquotedblleft{}if\textquotedblright{}. Today's \textquotedblleft{}if\textquotedblright{} is one that did not happen, and the second form comes back to say what would have followed.
\end{quote}

\subsection*{You'll want to know: an if that did not happen}
\begin{quote}
\textbf{\kn{ನನ್ನ} \kn{ಹತ್ತಿರ} \kn{ದುಡ್ಡು} \kn{ಇದ್ದಿದ್ದರೆ} \kn{ನಾನು} \kn{ಹೋಗುತ್ತಿದ್ದೆ}.}
\end{quote}

\begin{quote}
\emph{nanna hattira duḍḍu iddiddare nānu hōguttidde.} — \textquotedblleft{}If I had had money, I would have gone.\textquotedblright{}
\end{quote}

\textbf{\kn{ನನ್ನ} \kn{ಹತ್ತಿರ}}, \textquotedblleft{}near me\textquotedblright{}, is how Kannada says what you have on you: \textbf{\kn{ನನ್ನ} \kn{ಹತ್ತಿರ} \kn{ದುಡ್ಡು} \kn{ಇದೆ}}, \textquotedblleft{}I have money (on me)\textquotedblright{}.

For an \textquotedblleft{}if\textquotedblright{} that did not happen, Kannada puts \textbf{\kn{ಇದ್ದರೆ}} after the \textbf{-\kn{ಇ}} or \textbf{-\kn{ಉ}} form: \textbf{\kn{ಬಂದು}} + \textbf{\kn{ಇದ್ದರೆ}} $\to$ \textbf{\kn{ಬಂದಿದ್ದರೆ}} (\emph{bandiddare}, \textquotedblleft{}if (you) had come\textquotedblright{}); \textbf{\kn{ಮಾಡಿ}} + \textbf{\kn{ಇದ್ದರೆ}} $\to$ \textbf{\kn{ಮಾಡಿದ್ದರೆ}} (\emph{māḍiddare}, \textquotedblleft{}if (you) had done\textquotedblright{}); \textbf{\kn{ಇದ್ದು}} + \textbf{\kn{ಇದ್ದರೆ}} $\to$ \textbf{\kn{ಇದ್ದಿದ್ದರೆ}} (\emph{iddiddare}, \textquotedblleft{}if (there) had been\textquotedblright{}).

\subsection*{You'll want to know: would have}
The result takes \textbf{-\kn{ಉತ್ತಿದ್ದೆ}}, the form of \textquotedblleft{}I was doing\textquotedblright{} and \textquotedblleft{}I used to do\textquotedblright{}, here \textquotedblleft{}I would have done\textquotedblright{}: \textbf{\kn{ಹೋಗುತ್ತಿದ್ದೆ}}, \textquotedblleft{}I would have gone\textquotedblright{}. For \textquotedblleft{}it\textquotedblright{}, \textbf{-\kn{ಉತ್ತಿತ್ತು}}:

\begin{quote}
\textbf{\kn{ನೀನು} \kn{ಬಂದಿದ್ದರೆ} \kn{ನನಗೆ} \kn{ಖುಷಿ} \kn{ಆಗುತ್ತಿತ್ತು}.}
\end{quote}

\begin{quote}
\emph{nīnu bandiddare nanage khuṣi āguttittu.} — \textquotedblleft{}If you had come, I would have been glad.\textquotedblright{}
\end{quote}

The if-clause tells you which meaning you have. With \textbf{-\kn{ಇದ್ದರೆ}} in front of it, \textbf{\kn{ಹೋಗುತ್ತಿದ್ದೆ}} is \textquotedblleft{}would have gone\textquotedblright{}, not \textquotedblleft{}was going\textquotedblright{}.

\subsection*{Your turn}
Say these aloud:

\begin{itemize}
\raggedright
  \item \textquotedblleft{}If I had had money, I would have gone\textquotedblright{} — \emph{nanna hattira duḍḍu iddiddare nānu hōguttidde}
  \item \textquotedblleft{}If you had come, I would have been glad\textquotedblright{} — \emph{nīnu bandiddare nanage khuṣi āguttittu}
  \item \textquotedblleft{}If the bus had come, I would have come early\textquotedblright{} — \emph{bas bandiddare nānu bēga baruttidde}
\end{itemize}

\begin{tcolorbox}[breakable,colback=teal!4,colframe=teal!35!black,title={Before you move on}]
\textquotedblleft{}If (you) had come\textquotedblright{}? (\textbf{\kn{ಬಂದಿದ್ದರೆ}}.) \textquotedblleft{}If there had been\textquotedblright{}? (\textbf{\kn{ಇದ್ದಿದ್ದರೆ}}.) \textquotedblleft{}I would have gone\textquotedblright{}? (\textbf{\kn{ಹೋಗುತ್ತಿದ್ದೆ}}.) \textquotedblleft{}It would have been\textquotedblright{}? (\textbf{-\kn{ಉತ್ತಿತ್ತು}}: \textbf{\kn{ಆಗುತ್ತಿತ್ತು}}.) What tells you it means \textquotedblleft{}would\textquotedblright{}? (\textbf{The -\kn{ಇದ್ದರೆ} clause in front of it.})
\end{tcolorbox}
\section[(six lines)]{Six lines about a trip}

\label{lesson:KA-W318-six-lines}



\begin{quote}
The last connected message you wrote was about a familiar day. This one tells a story: six lines about a trip, to your home town or anywhere else.

Pick one reader: a friend, by name.
\end{quote}

\subsection*{What to know first}
A told story is not a list of events. It has a \textbf{scene}, with \textbf{-\kn{ಉವಾಗ}} or \textbf{-\kn{ಉತ್ತಿದ್ದೆ}}. It has \textbf{events} in the plain past, kept in order with \textbf{\kn{ಮೊದಲು}}, \textbf{\kn{ಅದರ} \kn{ನಂತರ}}, the \textbf{-\kn{ಇ}} / \textbf{-\kn{ಉ}} form and \textbf{\kn{ಕೊನೆಗೆ}}. It has at least one \textbf{reason}, with \textbf{\kn{ಹಾಗಾಗಿ}} or \textbf{\kn{ಯಾಕೆಂದರೆ}}, and a \textbf{feeling}, with \textbf{\kn{ನನಗೆ} … \kn{ಆಯಿತು}}.

A good one ends by looking at what might have been: \textbf{… -\kn{ಇದ್ದರೆ}, … -\kn{ಉತ್ತಿದ್ದೆ}}.

\subsection*{Writing — connected composition}
On paper, in Kannada script, write \textbf{six lines} to your reader about a trip, real or invented, one line for each of these:

\begin{enumerate}
\raggedright
  \item the scene: when it was, and what you were doing as it began, with \textbf{-\kn{ಉವಾಗ}};
  \item the first thing you did, with \textbf{\kn{ಮೊದಲು}}: what you cleaned or packed, and a ticket or a seat;
  \item what happened next, with \textbf{\kn{ಅದರ} \kn{ನಂತರ}} or an \textbf{-\kn{ಇ}} / \textbf{-\kn{ಉ}} form, and one thing that went wrong: a missed bus, a lost bag, or a fever and a medicine;
  \item why, or what followed, with \textbf{\kn{ಯಾಕೆಂದರೆ}} or \textbf{\kn{ಹಾಗಾಗಿ}};
  \item how it ended and how you felt, with \textbf{\kn{ಕೊನೆಗೆ}} and \textbf{\kn{ನನಗೆ} … \kn{ಆಯಿತು}};
  \item one line about what might have been, with \textbf{-\kn{ಇದ್ದರೆ} … -\kn{ಉತ್ತಿದ್ದೆ}}.
\end{enumerate}

There is no model on this page, and no word bank. Use only words you already write. The whole exercise is choosing the right past for each line, and joining the lines.

\begin{tcolorbox}[breakable,colback=teal!4,colframe=teal!35!black,title={Before you move on}]
Now read your lines aloud. Mark every verb \textbf{s} for scene and \textbf{e} for event. Line one should be mostly \textbf{s}; lines two and three mostly \textbf{e}.

Then two spelling checks. If you wrote \textbf{\kn{ಸ್ವಚ್ಛ}}, is the \textbf{\kn{ಛ}} small, under \textbf{\kn{ಚ}}? If a word opens with \emph{au}, as \textbf{\kn{ಔಷಧ}} does, it opens with the letter \textbf{\kn{ಔ}}; inside a word, \emph{au} is the sign \textbf{\kn{◌ೌ}}.

Last, read line six alone. Does the if-clause end in \textbf{-\kn{ಇದ್ದರೆ}}, and the main verb in \textbf{-\kn{ಉತ್ತಿದ್ದೆ}}?
\end{tcolorbox}
\section[(review)]{Review — it depends, and what would have been}

\label{lesson:KA-R318-depends-and-would}



\begin{quote}
Say \textquotedblleft{}to depend\textquotedblright{} (\emph{avalambisu}) and \textquotedblleft{}It depends on the weather\textquotedblright{}. (\emph{adu havāmānada mēle avalambiside}.)

\begin{itemize}
\raggedright
  \item \emph{Write it:} \textbf{\kn{ಛ}} and \textbf{\kn{ಔ}}, once each
\end{itemize}
\end{quote}

\subsection*{Your turn}
Say these aloud:

\begin{itemize}
\raggedright
  \item Will you come tomorrow? It depends on the rain. If it rains from morning until evening, I won't come.
  \item According to the timetable, the bus leaves at six. If you don't come quickly, it will go.
  \item If I had had time, I would have come.
  \item If you had come, we would have gone to the cinema together.
  \item read your six lines aloud to your reader
\end{itemize}

\begin{tcolorbox}[breakable,colback=teal!4,colframe=teal!35!black,title={Before you move on}]
\textquotedblleft{}It depends on …\textquotedblright{}? (\textbf{… }\kn{ಮೇಲೆ} \kn{ಅವಲಂಬಿಸಿದೆ}\textbf{.}) \textquotedblleft{}If (you) had come\textquotedblright{}? (\textbf{\kn{ಬಂದಿದ್ದರೆ}}.) \textquotedblleft{}I would have come\textquotedblright{}? (\textbf{\kn{ಬರುತ್ತಿದ್ದೆ}}.) \textquotedblleft{}According to\textquotedblright{}? (\textbf{\kn{ಪ್ರಕಾರ}}.)
\end{tcolorbox}
